\documentclass[11pt]{article}
\usepackage[margin=1in]{geometry}
\usepackage{amsmath}
\usepackage{amssymb}
\usepackage{amsthm}
\usepackage{hyperref}
\usepackage{listings}

\lstset{basicstyle=\ttfamily\footnotesize, breaklines=true, frame=single,
  columns=fullflexible, keepspaces=true, showstringspaces=false}

\theoremstyle{definition}
\newtheorem{definition}{Definition}
\theoremstyle{plain}
\newtheorem{theorem}{Theorem}

\title{Conjecture 00000002118 is false: \\
an exact refutation of the third-cumulant law for Quicksort comparisons}
\author{Zhiyuan Zhang (GitHub: LibertychaserUS)}
\date{2026-10-27}

\begin{document}
\maketitle

\noindent\textbf{Verdict.} The conjecture is \emph{refuted}. No constant $C$ makes
$\kappa_3(n)=2n-6\log n+C$ hold for all $n\ge 2$, for the natural logarithm or for any base $b>0$, $b\ne 1$.
The refutation is exact (rational arithmetic plus three elementary logarithm inequalities) and is
machine-checked in Lean~4 with Mathlib (v4.33.1). Only the axioms \texttt{propext}, \texttt{Classical.choice}
and \texttt{Quot.sound} are used.

\section{The conjecture as filed}

The filed statement (English part, quoted verbatim):

\begin{quote}
Definition: The Quicksort cost $Q_n$ is the number of comparisons. Conjecture: The third cumulant is
$\kappa_3(n) = 2n - 6\log n + C$ (with explicit constant $C$); the linear main coefficient of $\kappa_3$ is $2$.
\end{quote}

\section{Precise reading}

\paragraph{Cost.} $Q_n$ is written $C_n$ below. It is the number of comparisons of Quicksort with a uniformly random
pivot on $n$ distinct keys. Its law is the unique solution of
\[
  C_0 = C_1 = 0, \qquad C_n \overset{d}{=} (n-1) + C_J + C'_{\,n-1-J}\quad (n\ge 2),
\]
where $J$ is uniform on $\{0,\dots,n-1\}$, and $C$, $C'$ are independent copies of the cost. This is the standard
model (Knuth; R\"osler 1991; Neininger--R\"uschendorf 2004).

\paragraph{Third cumulant.} $\kappa_3(n)$ is the third cumulant of $C_n$, which equals its third central moment
$\mathbb{E}\big[(C_n-\mathbb{E}C_n)^3\big]$.

\paragraph{Logarithm.} The main reading uses the natural logarithm. The refutation also covers every base $b>0$, $b\ne 1$,
which settles the reading of $\log$ as $\log_2$ or $\log_{10}$ as well.

\paragraph{Claim refuted.} For the exact reading,
\[
  (\mathrm{R}1)\qquad \exists\, C\in\mathbb{R}\ \ \forall n\ge 2:\quad \kappa_3(n)=2n-6\log n+C .
\]
(R1) is false in the natural-log case, and (R2) below is false for every base. The conjecture as literally filed is
therefore false. The asymptotic reading ``$\kappa_3(n)=2n-6\log n+C+o(1)$'' is not refuted by a finite computation here;
Section~3 gives numerical evidence against it, but no formal proof.

\section{The counterexample and its verification}

\paragraph{Exact small values.} Solving the recursion in exact rational arithmetic gives:
\begin{center}
\begin{tabular}{c l c c}
\hline
$n$ & law of $C_n$ & $\mathbb{E}C_n$ & $\kappa_3(n)$ \\
\hline
2 & $1$ with probability $1$ & $1$ & $0$ \\
3 & $2$ w.p.\ $1/3$, $3$ w.p.\ $2/3$ & $8/3$ & $-2/27$ \\
4 & $4$ w.p.\ $1/2$, $5$ w.p.\ $1/6$, $6$ w.p.\ $1/3$ & $29/6$ & $13/54$ \\
5 & (not needed) & $-$ & $216/125 = 1.728$ \\
12 & (not needed) & $-$ & $136.954137045\ldots$ \\
\hline
\end{tabular}
\end{center}

\paragraph{Hand check of $n=2,3,4$.}
\begin{itemize}
\item $n=2$: $C_2=1+C_J+C'_{1-J}=1$ for both $J$.
\item $n=3$: $C_3=2+C_J+C'_{2-J}$ gives $J=1\Rightarrow 2$ (probability $1/3$), and $J\in\{0,2\}\Rightarrow 3$
(probability $2/3$). Central moment: deviations $-2/3$ (prob.\ $1/3$) and $+1/3$ (prob.\ $2/3$), so
$\kappa_3(3)=\tfrac13\cdot(-\tfrac{8}{27})+\tfrac23\cdot\tfrac1{27}=-\tfrac{2}{27}$.
\item $n=4$: $C_4=3+C_J+C'_{3-J}$. $J\in\{1,2\}$ gives $4$ (probability $1/2$); $J\in\{0,3\}$ gives $3+C_3$ (probability $1/2$),
so $5$ with probability $1/6$ and $6$ with probability $1/3$. Mean $29/6$, and
$\kappa_3(4)=\tfrac12(-\tfrac{125}{216})+\tfrac16\cdot\tfrac1{216}+\tfrac13\cdot\tfrac{343}{216}=\tfrac{13}{54}$.
\end{itemize}

\paragraph{Natural logarithm: contradiction at $n=2,3$.} The claimed identity at $n=2$ and $n=3$ reads
\[
  0 = 4-6\log 2+C,\qquad -\tfrac{2}{27} = 6-6\log 3+C .
\]
Subtracting gives $\log 3-\log 2=\tfrac{28}{81}\approx 0.345679$. But $\log(3/2)\approx 0.405465$, and the alternating
series bound for $\log(1-x)$ at $x=1/3$ with four terms gives $\log(3/2)\ge 0.4043-\tfrac1{162}>0.398$. So no real $C$ works.
Equivalently, the constant forced by $n=2$ is $6\log 2-4=0.158883$, and the one forced by $n=3$ is
$-\tfrac{2}{27}-6+6\log 3=0.517600$. They differ by $0.358717$.

\paragraph{Any base $b$: contradiction at $n=2,3,4$.} Write $L_n=\log n/\log b$. The identity at $n=2,3,4$, together with
$\log 4=2\log 2$, gives
\[
  0=4-6L_2+C,\quad -\tfrac{2}{27}=6-6L_3+C,\quad \tfrac{13}{54}=8-12L_2+C,
\]
so $L_2=\tfrac{203}{324}$ and $L_3=\tfrac{315}{324}$. Hence $203\log 3=315\log 2$, that is
$\log 3-\log 2=\tfrac{112}{203}\log 2$. Using $\log 2<0.6931471808$ (Mathlib: \texttt{Real.log\_two\_lt\_d9}) gives
$\log 3-\log 2<0.38249$, while $\log 3-\log 2\ge 0.398$ from the bound above. This is a contradiction. The base-$2$ and
base-$10$ cases are included. Numerically, $\log_2(3/2)=0.584963$ and $\log_{10}(3/2)=0.176091$, each incompatible with the
forced value $28/81$.

\paragraph{Growth (numerical evidence against the asymptotic reading).} Exact values give
\begin{center}
\begin{tabular}{r r r r}
\hline
$n$ & $\kappa_3(n)$ & $\kappa_3(n)/n^3$ & $\kappa_3(n)-(2n-6\log n)$ \\
\hline
12 & $136.954$ & $0.0793$ & $127.86$ \\
60 & $39877.35$ & $0.18462$ & $39781.9$ \\
90 & & $0.19891$ & \\
120 & $357056.64$ & $0.20663$ & $356845.4$ \\
\hline
\end{tabular}
\end{center}
Also $\kappa_3(120)/\kappa_3(60)=8.954$, which matches cubic growth (ratio $2^3=8$) and not linear growth (ratio $2$).
This is numerical evidence only; no proof of cubic growth is given here. (It is consistent with the known fact that
$(C_n-\mathbb{E}C_n)/n$ converges in law with convergent moments, and with the limit having non-zero third cumulant. We did
not verify that last point.)

\paragraph{Robustness to the cost convention.} Recomputing with toll $n+1$ (Knuth-style) or toll $n$ instead of $n-1$
changes the exact values, but not the conclusion. The constant forced by $n=2$ and by $n=3$ (natural log) is
\begin{center}
\begin{tabular}{l c c c}
\hline
toll & $\kappa_3(3)$ & $C$ from $n=2$ & $C$ from $n=3$ \\
\hline
$n-1$ (main) & $-2/27$ & $0.158883$ & $0.517600$ \\
$n+1$ & $-2$ & $0.158883$ & $-1.408326$ \\
$n$ & $-16/27$ & $0.158883$ & $-0.000919$ \\
\hline
\end{tabular}
\end{center}
In each case the two constants differ, so the refutation does not depend on the convention.

\section{Lean formalization}

The file is \texttt{Disproof.lean} in the batch directory. It uses \texttt{namespace TLMC2118}. The ASCII transcription
below is for readability; the Lean file itself is the authoritative text.

\subsection*{Definitions and correspondence}
\begin{lstlisting}
def qsPmf (n : Nat) : Nat -> Rat :=
  if h : n <= 1 then fun v => if v = 0 then 1 else 0
  else fun v => (1 / n) *
    sum j : Fin n, sum x in range (j*j+1),
      sum y in range ((n-1-j)*(n-1-j)+1),
        (if v = (n-1) + x + y then qsPmf j x * qsPmf (n-1-j) y else 0)
-- termination_by n  (decreasing_by discharged by omega)

def mom (n k : Nat) : Rat := sum v in range (n*n+1), (v : Rat)^k * qsPmf n v
def kappa3 (n : Nat) : Rat := sum v in range (n*n+1), ((v : Rat) - mom n 1)^3 * qsPmf n v
\end{lstlisting}

\begin{itemize}
\item \texttt{qsPmf n v} is $P(C_n=v)$. The base cases $n\le 1$ give $\delta_0$. The step for $n\ge 2$ is the convolution
$C_n=(n-1)+C_J+C'_{n-1-J}$ with $J$ uniform on $\texttt{Fin n}$, weight $1/n$, and independent copies.
\item The sums over $x$ and $y$ are truncated at $j^2$ and $(n-1-j)^2$. This loses no mass, because $C_j\le j(j-1)/2\le j^2$.
\item \texttt{kappa3 n} is the third central moment of $C_n$, computed exactly over \texttt{Rat}.
\end{itemize}

\subsection*{Computed values (Lean lemmas)}
\begin{lstlisting}
lemma qsPmf_two   : qsPmf 2 = fun v => if v = 1 then 1 else 0
lemma qsPmf_three : qsPmf 3 = fun v => if v = 2 then 1/3 else if v = 3 then 2/3 else 0
lemma qsPmf_four  : qsPmf 4 = fun v => if v = 4 then 1/2 else if v = 5 then 1/6
                                       else if v = 6 then 1/3 else 0
lemma kappa3_two   : kappa3 2 = 0
lemma kappa3_three : kappa3 3 = -2/27
lemma kappa3_four  : kappa3 4 = 13/54
lemma log_three_sub_log_two_gt : (0.39 : Real) < Real.log 3 - Real.log 2
\end{lstlisting}

\subsection*{Main theorems}
\begin{lstlisting}
theorem conjecture_00000002118_fails :
    not (exists C : Real, forall n : Nat, 2 <= n ->
      (kappa3 n : Real) = 2 * (n : Real) - 6 * Real.log (n : Real) + C)

theorem conjecture_00000002118_fails_any_base :
    not (exists b C : Real, 0 < b /\ b != 1 /\ forall n : Nat, 2 <= n ->
      (kappa3 n : Real) = 2 * (n : Real) - 6 * (Real.log (n : Real) / Real.log b) + C)
\end{lstlisting}

\noindent The first theorem is the negation of the conjecture (natural log) at the counterexample. The second is the same
for every base. Both are proved from the lemmas above and the logarithm bounds
(\texttt{Real.abs\_log\_sub\_add\_sum\_range\_le} for the lower bound on $\log(3/2)$, and \texttt{Real.log\_two\_lt\_d9}).
The proof uses no \texttt{decide} over large ranges and no \texttt{native\_decide}. Forbidden tokens do not occur.

\subsection*{Axioms}
\begin{lstlisting}
'TLMC2118.conjecture_00000002118_fails' depends on axioms: [propext, Classical.choice, Quot.sound]
'TLMC2118.conjecture_00000002118_fails_any_base' depends on axioms: [propext, Classical.choice, Quot.sound]
\end{lstlisting}

\noindent The file compiles with exit status 0 under \texttt{lean-toolchain} \texttt{v4.33.1}. The only message is a
harmless linter note about a \texttt{<;>} combinator.

\section{Reproduction}

\begin{enumerate}
\item Exact computation and robustness (Python 3, Fractions only):
\begin{lstlisting}
cd /Users/nihao/Projects/tlmc-work/batch/00000002118
/Users/nihao/Projects/tlmc-work/venv/bin/python check.py
/Users/nihao/Projects/tlmc-work/venv/bin/python convention_check.py
\end{lstlisting}
\item Lean verification:
\begin{lstlisting}
cd /Users/nihao/Projects/tlmc-work/template
lake env lean /Users/nihao/Projects/tlmc-work/batch/00000002118/Disproof.lean
\end{lstlisting}
\item This report:
\begin{lstlisting}
cd /Users/nihao/Projects/tlmc-work/batch/00000002118 && tectonic report.tex
\end{lstlisting}
\end{enumerate}

\section*{Limitations}
The formal result refutes the exact-identity reading of the conjecture, for all $n\ge 2$, with natural log or any base.
An asymptotic reading (``$\kappa_3(n)=2n-6\log n+C+o(1)$'') is not formally refuted here. The numerical growth evidence
is in Section~3. A rigorous proof of cubic growth would use the Neininger--R\"uschendorf convergence of moments, plus
the non-vanishing of the third cumulant of the limit law, which we did not verify.

\end{document}
